\conclusion

Цель выпускной квалификационной работы достигнута. В ходе работы разработана и реализована архитектурная основа распределённой платформы уведомлений, обеспечивающей устойчивую обработку уведомительных событий, контролируемую канальную доставку и прослеживаемость результатов.

Получены следующие результаты:
\begin{enumerate}
    \item определена доменная модель уведомительного процесса: уведомительное событие, запуск доставки, канальная доставка и история;
    \item сформулированы границы гарантированной доставки внутри платформы: сохранность принятого события, безопасная повторная обработка и фиксация результата;
    \item обоснована мультиканальная модель обработки на примере email- и SMS-контуров;
    \item спроектирована архитектура с разделением gRPC-командного контура и Kafka-контура фоновой обработки;
    \item реализованы сервисы notification-facade, template-registry, profile-consent, scheduler-delivery, notification-mail-sender, notification-sms-sender и history-writer;
    \item реализовано разделение хранилищ по сервисным контурам: отдельные PostgreSQL-хранилища для операционных сервисов, MongoDB для шаблонов и Redis для профилей получателей;
    \item реализованы механизмы надёжной обработки: transactional outbox, transactional inbox, idempotency, deduplication, retry и история статусов;
    \item проведена стендовая и нагрузочная проверка основных сценариев обработки.
\end{enumerate}

Проверка подтвердила, что реализованный контур обеспечивает прохождение полного маршрута обработки от входной команды до записи статуса в истории. Нагрузочный прогон показал наблюдаемость командного, событийного и канального участков обработки по метрикам gRPC, Kafka и JVM. По сохранённым артефактам зафиксированы активная фаза прогона длительностью около 44 минут, среднее потребление facade.events на уровне 38.3 req/s, среднее потребление delivery-statuses на уровне 21.5 req/s, средняя медианная задержка CreateNotificationEvent 786 ms и среднее время Kafka listener-обработки около 47 ms для sender- и history-контуров.

Границы результата соответствуют постановке работы: в текущей реализации поддержаны email и SMS; физическая доставка конечному пользователю зависит от внешнего провайдера; production-контур безопасности, промышленное масштабирование и подключение дополнительных каналов относятся к дальнейшему развитию.

Дальнейшее развитие платформы представляется целесообразным по следующим направлениям:
\begin{enumerate}
    \item подключение дополнительных каналов доставки без изменения базовой модели уведомительного события и истории;
    \item расширение набора количественных эксплуатационных метрик, сохраняемых вместе с артефактами нагрузочного прогона;
    \item развитие production-контуров безопасности и масштабирования.
\end{enumerate}
